% Packages and macros of the manual. Only packages from TeX Live's latex-base, latex-recommended,
% latex-extra, fonts-recommended and pictures collections (see README.md, "Manual").
\usepackage[utf8]{inputenc}
\usepackage[T1]{fontenc}
\usepackage{lmodern}
\usepackage{textcomp}
\usepackage{amssymb}
\usepackage[a4paper,margin=2.3cm,headheight=14pt]{geometry}
\usepackage{graphicx}
\usepackage[table]{xcolor}
\usepackage{microtype}
\usepackage{booktabs}
\usepackage{tabularx}
\usepackage{longtable}
\usepackage{array}
\usepackage[shortlabels]{enumitem}
\usepackage{caption}
\usepackage{listings}
\usepackage{tikz}
\usetikzlibrary{arrows.meta,positioning,fit,calc}
\usepackage[most]{tcolorbox}
\usepackage{fancyhdr}
\usepackage{makeidx}
\usepackage{url}
\usepackage[colorlinks=true,linkcolor=agernlink,urlcolor=agernlink,citecolor=agernlink,
  bookmarksnumbered=true,pdfusetitle=false]{hyperref}
\usepackage{bookmark}

\makeindex

% ---- colours (the app's amber accent on paper)
\definecolor{agernamber}{HTML}{C8761A}
\definecolor{agernlink}{HTML}{1F5F99}
\definecolor{agernink}{HTML}{1F2421}
\definecolor{notebg}{HTML}{EEF4FA}
\definecolor{warnbg}{HTML}{FDF1E3}
\definecolor{codebg}{HTML}{F5F5F2}

\hypersetup{
  pdftitle={Agern: Getting Started},
  pdfauthor={Egernet},
  pdfsubject={Getting-started manual for Agern, version \agernversion},
  pdfkeywords={Agern, software factory, Claude, GitHub},
}

% ---- layout
\setlength{\parindent}{0pt}
\setlength{\parskip}{0.55em plus 0.2em}
\setlength{\emergencystretch}{3em}
\setlist{itemsep=0.25em,topsep=0.4em}
\setlist[enumerate,1]{label=\textbf{\arabic*.},leftmargin=2em}
\captionsetup{font=small,labelfont=bf,margin=1em}
\captionsetup[figure]{skip=6pt}
\renewcommand{\arraystretch}{1.25}
\pagestyle{fancy}
\fancyhf{}
\fancyhead[L]{\small\nouppercase{\leftmark}}
\fancyhead[R]{\small Agern \agernversion}
\fancyfoot[C]{\small\thepage}
\renewcommand{\headrulewidth}{0.3pt}
\fancypagestyle{plain}{\fancyhf{}\fancyfoot[C]{\small\thepage}\renewcommand{\headrulewidth}{0pt}}

% ---- screenshots: every UI image of the manual goes through these two macros, one call per line,
% label last (scripts/manual.sh checks that each label is referenced and each file was produced):
%   \screenshot[<width>]{<file without .png>}{<caption>}{<label>}   -> figure with \label{fig:<label>}
%   \phoneshot[<width>]{<file without .png>}{<caption>}{<label>}    -> the same, narrow (390 px phone shots)
% The files come from build/manual/screenshots/ (web/scripts/screenshots.mjs).
\graphicspath{{../../build/manual/screenshots/}}
\newcommand{\screenshot}[4][\linewidth]{%
  \begin{figure}[htbp]
    \centering
    \setlength{\fboxsep}{0pt}\setlength{\fboxrule}{0.4pt}%
    \fcolorbox{black!30}{white}{\includegraphics[width=\dimexpr#1-0.8pt\relax,height=0.78\textheight,keepaspectratio]{#2}}
    \caption{#3}
    \label{fig:#4}
  \end{figure}}
\newcommand{\phoneshot}[4][0.40\linewidth]{%
  \begin{figure}[htbp]
    \centering
    \setlength{\fboxsep}{0pt}\setlength{\fboxrule}{0.4pt}%
    \fcolorbox{black!30}{white}{\includegraphics[width=\dimexpr#1-0.8pt\relax,height=0.78\textheight,keepaspectratio]{#2}}
    \caption{#3}
    \label{fig:#4}
  \end{figure}}

% ---- text markup
% \term{run}: a glossary term where it is introduced (italic, indexed). \term[index entry]{text}.
\newcommand{\term}[2][]{\emph{#2}\if\relax\detokenize{#1}\relax\index{#2}\else\index{#1}\fi}
% \ui{Start run}: a label of a button, field or page in the web UI.
\newcommand{\ui}[1]{\textsf{\textbf{#1}}}
% \cmd{sudo agern doctor}, \envvar{FACTORY_PUBLIC_URL}: commands and variables (underscores allowed).
\newcommand{\cmd}[1]{\texttt{\detokenize{#1}}}
% (\ttraw is robust, so the index file gets \ttraw{NAME} rather than an expanded name.)
\DeclareRobustCommand{\ttraw}[1]{\texttt{\detokenize{#1}}}
\newcommand{\envvar}[1]{\ttraw{#1}\index{#1@\ttraw{#1}}}
\newcommand{\file}[1]{\path{#1}}
% \msg{...}: a message exactly as Agern shows it.
\newcommand{\msg}[1]{\textsf{``#1''}}
% Links into the repository's reference docs.
\newcommand{\repodoc}[2]{\href{https://github.com/Egernet-CS/agern/blob/main/docs/#1}{#2}}

% ---- boxes
\newtcolorbox{note}[1][Note]{enhanced,breakable,colback=notebg,colframe=agernlink,boxrule=0.4pt,
  left=6pt,right=6pt,top=4pt,bottom=4pt,fonttitle=\bfseries\small,title=#1,arc=2pt}
\newtcolorbox{warning}[1][Warning]{enhanced,breakable,colback=warnbg,colframe=agernamber,boxrule=0.6pt,
  left=6pt,right=6pt,top=4pt,bottom=4pt,fonttitle=\bfseries\small,title=#1,arc=2pt}

% ---- code
\lstdefinestyle{sh}{language=bash,basicstyle=\ttfamily\small,columns=fullflexible,keepspaces=true,
  upquote=true,breaklines=true,breakatwhitespace=false,showstringspaces=false,backgroundcolor=\color{codebg},
  frame=single,rulecolor=\color{black!15},framesep=4pt,xleftmargin=4pt,xrightmargin=4pt,
  commentstyle=\color{black!55},
  breakindent=2em,aboveskip=0.6em,belowskip=0.6em}
\lstdefinestyle{conf}{style=sh,language={}}
\lstset{style=sh}

% ---- "What you have now" ends every chapter.
\newcommand{\whatyouhave}{\section*{What you have now}\addcontentsline{toc}{section}{What you have now}}

% Characters pdfLaTeX's utf8 does not map by itself.
\DeclareUnicodeCharacter{2192}{\ensuremath{\rightarrow}}
\DeclareUnicodeCharacter{2026}{\ldots}
\DeclareUnicodeCharacter{00B7}{\textperiodcentered}
\DeclareUnicodeCharacter{2713}{\ensuremath{\checkmark}}
\DeclareUnicodeCharacter{25B6}{\ensuremath{\blacktriangleright}}
